\chapter{Thảo luận, mô hình triển khai và hướng phát triển}
\label{chap:thao_luan}

Chương này thảo luận sâu về ý nghĩa học thuật, tính kinh tế phần cứng trong triển khai thực tế, tính công bằng của các thực nghiệm đối sánh, đồng thời thẳng thắn chỉ ra các giới hạn hiện tại của công trình nghiên cứu và xác định các hướng mở rộng tiếp nối trong tương lai.

% ==============================================================================
\section{Phác thảo thiết kế kiến trúc phân tán Edge--Cloud (Định hướng tương lai)}
\label{sec:kien_truc_trien_khai_edge_cloud}

Để làm rõ động lực của các lựa chọn thiết kế được đánh giá trong thực nghiệm (tính kinh tế về tham số, bộ nhớ kích hoạt GPU và cơ chế truyền tải chỉ chứa vector số đếm), nghiên cứu phác thảo ý niệm về một kiến trúc phân tán giữa thiết bị biên và máy chủ trung tâm (Edge--Cloud Split). \textbf{Cần nhấn mạnh rằng đây là bản phác thảo thiết kế định hướng ứng dụng trong tương lai (Future Work), không phải kết quả đo đạc thực nghiệm trên phần cứng nhúng}. Toàn bộ các kết quả thực nghiệm được báo cáo trong đề tài đều được thực hiện trên hệ thống xử lý tập trung chuẩn hóa đã mô tả tại Chương \ref{chap:phuong_phap} và Chương \ref{chap:thuc_nghiem}.

\begin{figure}[htbp]
\centering
\includegraphics[width=0.98\linewidth]{fig/edge_cloud_architecture.pdf}
\caption{Bản phác thảo thiết kế kiến trúc phân tán Edge--Cloud định hướng ứng dụng: xử lý nhận dạng tại thiết bị biên và dự báo dòng giao thông trên đồ thị tại máy chủ trung tâm.}
\label{fig:edge_cloud_arch_chap6}
\end{figure}

Trong kiến trúc định hướng triển khai (Hình \ref{fig:edge_cloud_arch_chap6}):
\begin{itemize}
    \item \textbf{Tại thiết bị biên (Edge Devices -- các trạm camera quan sát):} Mô hình Giai đoạn 1 (mạng xương sống EfficientNet-B5 hoặc phương án siêu nhẹ MobileNetV3-Large sau khi lượng tử hóa số nguyên INT8 hoặc FP16 thông qua các runtime nhúng chuyên dụng như OpenVINO, TensorRT) được nạp vào bộ nhớ thiết bị biên. Thiết bị thực hiện suy luận ảnh chụp tĩnh trực tiếp trong bộ nhớ đệm và chỉ phát đi gói tin nhị phân chứa vector số lượng phương tiện theo phân lớp: $[\hat{c}_{\text{car}}, \hat{c}_{\text{bike}}]^T$.
    \item \textbf{Tại máy chủ trung tâm (Central Cloud Server):} Tiếp nhận các vector số đếm từ toàn bộ $N=608$ nút camera trên mạng lưới đường bộ đô thị, cập nhật vào ma trận trạng thái chuỗi thời gian và thực thi mô hình TA-STGCN để dự báo lưu lượng toàn mạng lưới trong 30 phút tiếp theo.
\end{itemize}

\subsection{Ý nghĩa đối với băng thông truyền thông hạ tầng}
Trong các giải pháp giám sát giao thông truyền thống, việc truyền tải liên tục luồng video độ nét cao hoặc ảnh tĩnh độ phân giải lớn từ hàng trăm điểm quan sát về trung tâm chỉ huy gây áp lực rất lớn lên hạ tầng mạng viễn thông. 

Bằng cách chuyển dịch khâu nhận dạng xuống thiết bị biên, mỗi nút quan sát chỉ cần truyền tải vector số đếm hai chiều thay vì truyền hình ảnh thô. Về mặt ý niệm lý thuyết, việc truyền một vector hai số thực (vài bytes) so với một ảnh nén JPEG kích thước hàng trăm kilobytes đem lại tiềm năng cắt giảm lưu lượng truyền dẫn rất lớn. Tuy nhiên, tỷ lệ cắt giảm lý thuyết này chưa tính đến các chi phí giao thức truyền thông (protocol overhead), cơ chế xác thực, gói tin điều khiển mạng và việc truyền lại khi rớt gói. Do đó, việc lượng hóa chính xác băng thông thực tế đòi hỏi các đo đạc thực nghiệm toàn diện trên môi trường mạng thực địa trong các nghiên cứu tương lai.

\subsection{Bảo vệ quyền riêng tư hình ảnh theo thiết kế (Privacy-by-Design)}
Việc truyền tải và lưu trữ tập trung hình ảnh camera nơi công cộng tiềm ẩn nguy cơ làm lộ lọt thông tin cá nhân của người tham gia giao thông (nhận diện khuôn mặt hoặc biển số phương tiện).

Trong kiến trúc phân tán được phác thảo, nguyên lý bảo vệ quyền riêng tư được bảo đảm thông qua cơ chế xử lý nội bộ tại chỗ:
\begin{itemize}
    \item Khung hình camera được giải mã trực tiếp trong bộ nhớ RAM tạm thời của vi xử lý biên, đưa qua mạng nơ-ron hồi quy đếm xe và được giải phóng ngay lập tức sau khi có kết quả.
    \item Không có bất kỳ hình ảnh thô hay đặc trưng nhận dạng cá nhân nào được lưu trữ cố định trên bộ nhớ biên hoặc truyền dẫn qua môi trường mạng Internet.
    \item Dữ liệu rời khỏi thiết bị biên chỉ là các con số thống kê tổng hợp về lưu lượng phương tiện của tuyến đường, hoàn toàn loại bỏ nguy cơ rò rỉ dữ liệu hình ảnh nhạy cảm.
\end{itemize}

% ==============================================================================
\section{Tính kinh tế phần cứng: Bộ nhớ kích hoạt GPU đỉnh so với số lượng tham số tĩnh}
\label{sec:tinh_kinh_te_phan_cung}

Một phát hiện thực nghiệm quan trọng từ các kết quả tại Chương \ref{chap:thuc_nghiem} (Bảng \ref{tab:sieu_tham_so_chap5_models}) là: \textbf{Số lượng tham số tĩnh ($N_{\text{params}}$) không phản ánh đầy đủ chi phí phần cứng thực tế của các mô hình đồ thị không -- thời gian trong môi trường đô thị quy mô lớn}.

\begin{enumerate}
    \item \textbf{Hiện tượng bùng nổ bộ nhớ kích hoạt (Activation Memory Explosion):} Các mô hình Transformer đồ thị tiên tiến như STAEformer duy trì lượng tham số tĩnh ở mức vừa phải ($568\text{K}$ tham số), song việc tính toán các ma trận chú ý mềm liên hợp (Joint Spatio-Temporal Attention) trên toàn bộ $N=608$ nút khiến lượng bộ nhớ kích hoạt GPU đỉnh vọt lên tới \textbf{11.509,1 MB}. Mức tiêu thụ này chạm sát ngưỡng dung lượng 12 GB của các dòng card đồ họa thương mại phổ biến, gây khó khăn lớn cho việc triển khai trên các cụm máy chủ kinh phí hạn chế.
    \item \textbf{Bản chất thiết kế của mô hình iTransformer:} Mặc dù iTransformer đạt mức tiêu thụ bộ nhớ GPU cực kỳ thấp (\textbf{265,2 MB}), điều này hoàn toàn xuất phát từ bản chất thiết kế Inverted Transformer của mô hình: coi mỗi nút giao thông như một biến số độc lập và áp dụng cơ chế chú ý trên chiều biến số, hoàn toàn không xây dựng hay duy trì ma trận không gian $N \times N$. Do không mô hình hóa tương quan không gian qua đồ thị, iTransformer là một mô hình chuỗi thời gian đa biến thuần túy (non-graph) và được xử lý riêng biệt trong các đánh giá.
    \item \textbf{Cân bằng tối ưu của TA-STGCN:} Bằng cách tách rời phép tích chập không gian Chebyshev ($K=3$) trên đồ thị thưa khỏi cơ chế Self-Attention trên chiều thời gian độc lập cho từng nút, TA-STGCN chỉ tiêu tốn \textbf{8.606,9 MB} bộ nhớ GPU đỉnh (tiết kiệm $25.2\%$ so với STAEformer) và giảm $70.4\%$ số lượng tham số so với ASTGCN ($453\text{K}$ so với $1.533\text{K}$). Sự cân đối đồng thời trên cả hai trục tham số và bộ nhớ kích hoạt giúp mô hình phù hợp với các hệ thống máy chủ vận hành thực tế.
\end{enumerate}

% ==============================================================================
\section{Tính khách quan và công bằng trong so sánh thực nghiệm}
\label{sec:tinh_cong_bang_thuc_nghiem}

Để đảm bảo tính nghiêm cẩn học thuật, các thực nghiệm đối sánh ở Giai đoạn 2 được thiết kế như một \textbf{thực nghiệm đối chứng có kiểm soát chặt chẽ (controlled implementation comparison)} thay vì một cuộc chạy đua tối ưu hóa riêng lẻ từng mô hình:

\begin{enumerate}
    \item \textbf{Sử dụng chung hàm mục tiêu (Unified Huber Objective):} Tất cả các mô hình đối chứng đều được huấn luyện dưới cùng một hàm mất mát Huber Loss ($\delta=1.0$). Một số mô hình gốc trong y văn ban đầu được đề xuất với hàm mục tiêu khác (chẳng hạn Graph WaveNet ban đầu sử dụng hàm mất mát $L_1$). Việc huấn luyện lại toàn bộ mô hình dưới hàm Huber Loss là một sự đánh đổi có chủ đích: kiểm soát đồng nhất cách mô hình phản ứng với các điểm ngoại lai để cô lập yếu tố sai khác về mặt kiến trúc. Do đó, kết quả của các mô hình đối chứng được hiểu là hiệu năng đạt được dưới hàm mục tiêu dùng chung, không phải kết quả công bố riêng lẻ của từng bài báo gốc.
    \item \textbf{Không tinh chỉnh siêu tham số riêng biệt cho từng baseline:} Nghiên cứu không thực hiện tìm kiếm siêu tham số vét cạn riêng lẻ cho từng mô hình đối chứng trên đồ thị 608 nút; các mô hình sử dụng cấu hình chuẩn do các tác giả nguyên bản đề xuất, được điều chỉnh kích thước đầu vào và đầu ra cho tương thích. Do đó, một số mô hình đối chứng có thể đạt độ chính xác cao hơn nếu được tinh chỉnh sâu. Đây là lý do nghiên cứu không nhấn mạnh vào các chênh lệch sai số nhỏ trên dữ liệu sạch, phù hợp với kết quả kiểm định kích thước ảnh hưởng ($|d| < 0.1$) cho thấy các mô hình đồ thị đạt độ chính xác điểm gần như tương đương nhau.
    \item \textbf{Phân định rõ ràng bản chất non-graph của iTransformer:} iTransformer được đưa vào thực nghiệm như một mô hình tham chiếu mạnh về chuỗi thời gian đa biến, không phải là mô hình dự báo cấu trúc đồ thị. Mô hình này được loại trừ khỏi các kết luận so sánh nội bộ giữa các kiến trúc đồ thị. Việc iTransformer đạt độ bền vững tốt trước nhiễu (ở Mức 4) cung cấp bằng chứng cho thấy một phần năng lực kháng nhiễu bắt nguồn từ cơ chế chú ý theo trục thời gian nói chung chứ không riêng gì cấu trúc đồ thị.
\end{enumerate}

% ==============================================================================
\section{Các giới hạn học thuật của nghiên cứu và Hướng phát triển tiếp nối}
\label{sec:han_che_va_huong_phat_trien}

Nghiên cứu thẳng thắn nhìn nhận các giới hạn kỹ thuật hiện tại và xác định các hướng nghiên cứu mở rộng trong tương lai:

\begin{enumerate}
    \item \textbf{Giới hạn trần sai số phụ thuộc giữa hai giai đoạn:} Mô hình dự báo Giai đoạn 2 được huấn luyện trực tiếp trên chuỗi dữ liệu số đếm sinh ra từ Giai đoạn 1. Do đó, độ chính xác và độ tin cậy của toàn bộ hệ thống dự báo bị giới hạn trên bởi biên độ sai số của khâu nhận dạng thị giác ban đầu.
    \item \textbf{Thách thức quy mô tính toán khi triển khai thực tế trên diện rộng:} Mặc dù mô hình đồ thị Giai đoạn 2 tương đối gọn nhẹ ($453\text{K}$ tham số), việc triển khai đồng loạt mô hình Giai đoạn 1 (EfficientNet-B5 với 28.3M tham số) trên 608 camera biên tương đương với việc vận hành tổng cộng hơn $17$ tỷ tham số trên toàn đô thị. Đây là một thách thức lớn về chi phí phần cứng biên, độ trễ và năng lượng tiêu thụ. Mặc dù MobileNetV3-Large (4.2M tham số) là một lựa chọn nhỏ gọn hơn, mô hình này phải đánh đổi bằng mức sai số Stage 1 cao hơn ($4.0126$ so với $3.6363$). Việc lượng tử hóa (INT8/FP16), chưng cất tri thức (Knowledge Distillation) và đo đạc thông lượng thực tế trên các bo mạch nhúng (như NVIDIA Jetson, Coral TPU) là hướng đi bắt buộc trước khi đưa vào ứng dụng thực tế.
    \item \textbf{Đánh đổi của phương pháp hồi quy đếm xe trực tiếp:} Nghiên cứu định vị cơ chế hồi quy đếm xe trực tiếp từ ảnh tĩnh như một giải pháp thực tế giúp loại bỏ chi phí gán nhãn điểm/bounding box phức tạp và không đòi hỏi căn chỉnh thông số hình học từng camera. Tuy nhiên, nghiên cứu chưa thực hiện đối sánh đối đầu trực tiếp với các phương pháp ước lượng bản đồ mật độ (Density-map estimators như MCNN \cite{mcnn}, CSRNet \cite{csrnet}, CAN \cite{can}) hoặc các mô hình phát hiện kết hợp bám vết (Tracking-by-Detection như YOLO \cite{yolo} + ByteTrack) trên cùng một tập dữ liệu gán nhãn vi mô.
    \item \textbf{Phạm vi không gian -- thời gian của bộ dữ liệu:} Bộ dữ liệu thực nghiệm được thu thập tại một đô thị đơn lẻ (TP. Hồ Chí Minh) trong khoảng thời gian 3 tháng, chưa bao quát đầy đủ các biến thiên thời tiết cực đoan theo mùa (mùa mưa ngập kéo dài hoặc sương mù dày đặc).
    \item \textbf{Các dòng kiến trúc chuỗi thời gian mới (2024--2025):} Thực nghiệm của đề tài bao quát các kiến trúc đồ thị và Transformer tiêu biểu đến năm 2024, song chưa tích hợp các dòng mô hình mới nhất xuất hiện trong giai đoạn 2024--2025 như Spatio-Temporal State-Space Models (Mamba \cite{mamba}), Diffusion STG, hoặc các mô hình chuỗi thời gian dựa trên Patch / LLM (như PatchTST \cite{patchtst}, TimeXer). Các họ kiến trúc này chưa được chuẩn hóa cho dữ liệu đồ thị có cấu trúc phân tách chủng loại xe thu thập từ camera, đòi hỏi những nghiên cứu kỹ thuật sâu hơn trong tương lai.
    \item \textbf{Tính thích ứng của ma trận kề không gian:} Đồ thị RBF xây dựng từ OpenStreetMap thể hiện cấu trúc kết nối hình học tĩnh, chưa thể tự thích ứng với các biến cố phi vật lý hoặc sự cố đột xuất (tai nạn giao thông làm chặn đường, ngập nước cục bộ). Hướng phát triển tiếp theo là nghiên cứu kết hợp đồ thị tĩnh với cơ chế học đồ thị thích ứng động (Dynamic Adaptive Graph Learning), đồng thời mở rộng phân tích độ bền vững dưới các phân phối nhiễu xác suất khác (như Weibull hoặc Log-normal).
\end{enumerate}
